\documentclass[12pt,a4paper]{article}

\usepackage[utf8]{inputenc}
\usepackage[T1]{fontenc}
\usepackage[spanish]{babel}
\usepackage{amsmath, amssymb}
\usepackage{geometry}
\usepackage{booktabs}
\usepackage{array}
\usepackage{xcolor}
\usepackage{hyperref}
\usepackage{enumitem}
\usepackage{fancyhdr}
\usepackage{graphicx}
\usepackage{tcolorbox}
\usepackage{multirow}
\usepackage{caption}
\usepackage{tikz}

\geometry{top=2.5cm, bottom=2.5cm, left=3cm, right=3cm}

% ── Colores ──────────────────────────────────────────────────────────────────
\definecolor{uddblue}{RGB}{0, 75, 135}
\definecolor{lightgray}{RGB}{245, 245, 245}

% ── Cajas de ayuda ────────────────────────────────────────────────────────────
\tcbuselibrary{skins, breakable}
\newtcolorbox{ayuda}[1]{
    colback   = lightgray,
    colframe  = uddblue,
    fonttitle = \bfseries\small,
    title     = {#1},
    breakable,
    left      = 6pt,
    right     = 6pt,
    top       = 4pt,
    bottom    = 4pt
}

% ── Diagrama de amenazas (tablero 5x5 con la pieza al centro) ────────────────
% #1 = lista de casillas amenazadas, #2 = letra de la pieza
\newcommand{\tablero}[2]{%
\begin{tikzpicture}[scale=0.40]
    \foreach \c/\r in {#1}{\fill[uddblue!25] (\c,\r) rectangle ++(1,1);}
    \fill[uddblue!70] (2,2) rectangle ++(1,1);
    \draw[step=1, gray!60] (0,0) grid (5,5);
    \draw[gray!60] (0,0) rectangle (5,5);
    \node[white, font=\bfseries\small] at (2.5,2.5) {#2};
\end{tikzpicture}}

% ── Encabezado y pie ─────────────────────────────────────────────────────────
\pagestyle{fancy}
\fancyhf{}
\fancyhead[L]{\small\color{gray} IIP314W\,--\,Optimización Aplicada a Negocios}
\fancyhead[R]{\small\color{gray} Tarea 3\,--\,2026-T2}
\fancyfoot[C]{\small\thepage}
\renewcommand{\headrulewidth}{0.4pt}

\hypersetup{colorlinks=true, linkcolor=uddblue, urlcolor=uddblue}

\setlength{\parskip}{0.5em}
\setlength{\parindent}{0pt}

% =============================================================================
\begin{document}
% =============================================================================

% ── Título ───────────────────────────────────────────────────────────────────
\begin{center}
    {\Large\bfseries\color{uddblue}
     (IIP314W) Optimización Aplicada a Negocios}\\[0.4em]
    {\large Universidad del Desarrollo}\\[0.8em]
    \noindent\rule{\linewidth}{1.5pt}\\[0.5em]
    {\LARGE\bfseries Tarea 3}\\[0.3em]
    {\large Ajedrez y Optimización Entera:\\
     Máxima Ocupación de un Tablero sin Amenazas}\\[0.5em]
    \noindent\rule{\linewidth}{1.5pt}\\[0.8em]
    \begin{tabular}{@{}ll@{}}
        \textbf{Profesor:} & Ing.\ Rodrigo Trigo Vilches \\[0.3em]
        \textbf{Ayudante:} & Lic.\ Vicente Ramírez Almonacid \\[0.3em]
        \textbf{Período:}  & Año 2026\,--\,Trimestre 2
    \end{tabular}
\end{center}

\vspace{1em}
\noindent\rule{\linewidth}{0.4pt}
\vspace{0.5em}

% =============================================================================
\section*{Contexto}
% =============================================================================

En esta tarea deberán \textbf{modelar y resolver problemas de optimización lineal
entera} para determinar la \textbf{máxima cantidad de piezas de ajedrez} que se
pueden ubicar en un tablero estándar de $8\times8$ de modo que \textbf{ninguna
amenace a otra}.

Aunque el enunciado esté planteado sobre un tablero, la estructura del problema es
exactamente la misma que aparece en múltiples problemas de negocio: se trata de
\textbf{seleccionar la mayor cantidad de elementos posible de un conjunto, sujeto a
que ciertos pares son incompatibles entre sí}. Es el mismo esqueleto de la
asignación de antenas que no pueden interferirse, de turnos que no pueden
solaparse, de productos químicos que no pueden almacenarse juntos o de máquinas que
no pueden operar simultáneamente. Lo que cambia es la historia; el modelo es el
mismo. Por eso el ajedrez es un excelente banco de pruebas: las reglas son
inequívocas y la solución se puede verificar a simple vista.

\begin{ayuda}{Observación: el bloqueo no cambia el problema}
\vspace{2pt}
Al modelar podría surgir la duda de si una pieza ``bloquea'' la línea de otra y la
protege. No hace falta preocuparse por eso: como \emph{todas} las piezas del tablero
son piezas que ustedes colocaron, el bloqueo es \textbf{irrelevante}. En efecto, una
pieza deslizante (torre, alfil o reina) siempre amenaza a la pieza \textbf{más
cercana} sobre su línea de acción; si hubiera una pieza intermedia, esa pieza
intermedia ya estaría amenazada, y eso por sí solo viola la condición del problema.

En consecuencia, exigir que \textbf{ninguna pieza amenace a otra} es
\textbf{equivalente} a exigir que \textbf{no haya dos piezas sobre una misma línea de
acción}, y pueden modelar directamente esta última condición.
\end{ayuda}

% =============================================================================
\section*{Recordatorio: movimiento y amenaza de las piezas}
% =============================================================================

En los diagramas, la casilla \textbf{oscura} es la posición de la pieza y las
casillas \textbf{claras} son las que esa pieza amenaza.

\begin{center}
\begin{tabular}{cccc}
\tablero{0/2,1/2,3/2,4/2,2/0,2/1,2/3,2/4}{T} &
\tablero{0/0,1/1,3/3,4/4,0/4,1/3,3/1,4/0}{A} &
\tablero{0/2,1/2,3/2,4/2,2/0,2/1,2/3,2/4,0/0,1/1,3/3,4/4,0/4,1/3,3/1,4/0}{R} &
\tablero{1/0,3/0,0/1,4/1,0/3,4/3,1/4,3/4}{C} \\[0.3em]
\textbf{Torre} & \textbf{Alfil} & \textbf{Reina} & \textbf{Caballo}
\end{tabular}
\end{center}

\begin{center}
\renewcommand{\arraystretch}{1.35}
\begin{tabular}{@{}lp{10.2cm}@{}}
\toprule
\textbf{Pieza} & \textbf{Cómo se mueve y amenaza} \\
\midrule
\textbf{Torre} & Se mueve en línea recta a lo largo de su \textbf{fila} y de su
\textbf{columna}. Amenaza toda casilla que comparta fila o columna con ella. \\
\textbf{Alfil} & Se mueve en \textbf{diagonal}. Amenaza toda casilla que esté sobre
alguna de sus dos diagonales. Un alfil nunca cambia el color de casilla. \\
\textbf{Reina} & Combina la torre y el alfil: se mueve por \textbf{fila, columna y
diagonales}. Es la pieza de mayor alcance. \\
\textbf{Caballo} & Se mueve en \textbf{``L''}: dos casillas en una dirección y una
en la perpendicular. Amenaza hasta 8 casillas y es la única pieza que
\textbf{salta}: su amenaza no depende de las casillas intermedias. \\
\bottomrule
\end{tabular}
\captionof{table}{Movimiento y amenaza de las cuatro piezas consideradas.}
\end{center}

\noindent\textit{Nota sobre coordenadas:} conviene indexar las casillas por
\textbf{fila} $i$ y \textbf{columna} $j$, con $i,j\in\{1,\dots,8\}$. Con esa
notación, dos casillas $(i,j)$ y $(k,l)$ están en la misma diagonal si
$i-j=k-l$ o bien $i+j=k+l$.

\noindent\rule{\linewidth}{0.4pt}

% =============================================================================
\section*{Parte A \quad Formulación de los Modelos \hfill\normalfont\normalsize(40 pts)}
% =============================================================================

Formule un \textbf{modelo de optimización lineal entera} para cada uno de los
siguientes cuatro casos, en un tablero de $8\times8$:

\begin{enumerate}[label=\alph*), leftmargin=*, topsep=2pt]
    \item \textbf{Torres:} máxima cantidad de torres sin que ninguna amenace a otra.
    \item \textbf{Alfiles:} máxima cantidad de alfiles sin amenazas mutuas.
    \item \textbf{Reinas:} máxima cantidad de reinas sin amenazas mutuas.
    \item \textbf{Caballos:} máxima cantidad de caballos sin amenazas mutuas.
\end{enumerate}

Para \textbf{cada} modelo, su formulación debe ser completamente algebraica (en
notación matemática, \textbf{no} en código) y debe contener:

\begin{enumerate}[label=\arabic*., leftmargin=*, topsep=2pt]
    \item \textbf{Conjuntos e índices:} cómo indexa las casillas del tablero.
    \item \textbf{Variables de decisión:} defina con precisión su significado y su
          \textbf{tipo} y \textbf{dominio}. Piense qué representa cada variable y
          por qué debe ser del tipo que usted propone.
    \item \textbf{Función objetivo:} la expresión a maximizar.
    \item \textbf{Restricciones:} escriba y \textbf{nombre} las restricciones que
          impiden que dos piezas se amenacen. Explique, para cada pieza, cuál es la
          \textbf{condición geométrica} que define una amenaza y cómo la traduce a
          desigualdades lineales.
\end{enumerate}

\begin{ayuda}{Sugerencia de método (no de solución)}
\vspace{2pt}
Note que las cuatro formulaciones comparten el mismo esqueleto y solo cambian en
\emph{qué pares de casillas son incompatibles}. Conviene definir primero, de forma
general, qué significa que ``la casilla $(i,j)$ y la casilla $(k,l)$ estén en
conflicto para la pieza $P$'' y recién entonces particularizar a cada pieza.
Reflexione además si conviene escribir una restricción \textbf{por cada par de
casillas en conflicto} o si, para alguna pieza en particular, existe una forma más
\textbf{compacta y agregada} de expresar lo mismo (por ejemplo, agrupando por fila,
columna o diagonal). Ambas alternativas son válidas; discutir la diferencia suma
puntos.
\end{ayuda}

\noindent\rule{\linewidth}{0.4pt}

% =============================================================================
\section*{Parte B \quad Resolución con GurobiPy \hfill\normalfont\normalsize(30 pts)}
% =============================================================================

Implemente y resuelva los cuatro modelos de la Parte~A con \texttt{gurobipy}.
Responda:

\begin{enumerate}[label=\alph*), leftmargin=*, topsep=2pt]
    \item Reporte, para cada pieza, la \textbf{cantidad máxima} obtenida.
          Preséntelo en una tabla comparativa.
    \item Muestre, para cada caso, \textbf{una configuración óptima} del tablero
          (una disposición concreta de las piezas que alcanza ese máximo).
    \item \textbf{Verifique} que las soluciones son válidas: compruebe
          explícitamente que ningún par de piezas se amenaza.
    \item Analice y explique los resultados: ¿por qué unas piezas admiten muchas
          más que otras? Ordene las piezas de menor a mayor cantidad máxima y
          relacione ese orden con el \textbf{alcance} de cada pieza.
    \item Para el caso de los \textbf{caballos}, observe la configuración óptima
          que obtuvo. ¿Nota algún patrón en el \textbf{color de las casillas}
          ocupadas? Explique a qué se debe.
    \item ¿La solución óptima es \textbf{única} en cada caso? Si encuentra que no
          lo es, explique qué significa que existan múltiples óptimos.
\end{enumerate}

\noindent\rule{\linewidth}{0.4pt}
\newpage

% =============================================================================
\section*{Parte C \quad Combinación de Piezas por Valor \hfill\normalfont\normalsize(30 pts)}
% =============================================================================

Hasta aquí cada tablero contenía \textbf{un solo tipo} de pieza. Ahora
generalizaremos: podremos colocar \textbf{piezas de distintos tipos
simultáneamente} sobre el mismo tablero, y ya no maximizaremos la \emph{cantidad}
de piezas sino el \textbf{valor total} del conjunto colocado.

Usaremos la valoración clásica del ajedrez:

\begin{center}
\renewcommand{\arraystretch}{1.2}
\begin{tabular}{lcccc}
\toprule
\textbf{Pieza} & Reina & Torre & Alfil & Caballo \\
\midrule
\textbf{Valor} & 9 & 5 & 3 & 3 \\
\bottomrule
\end{tabular}
\captionof{table}{Valor relativo de cada pieza.}
\end{center}

Las reglas se mantienen: en cada casilla puede haber \textbf{a lo más una pieza}, y
\textbf{ninguna pieza puede amenazar a otra}, sin importar de qué tipo sea cada una.

\begin{ayuda}{Ojo con las amenazas entre piezas distintas}
\vspace{2pt}
Cuando conviven piezas de distintos tipos, la condición de conflicto deja de ser
simétrica en un solo sentido: dos piezas en las casillas $(i,j)$ y $(k,l)$ son
incompatibles si \textbf{la primera amenaza a la segunda \emph{o} la segunda amenaza
a la primera}. Por ejemplo, una torre y un caballo en la misma fila se amenazan
(por la torre), aunque el caballo no ataque a la torre. Su modelo debe capturar
ese ``o''.
\end{ayuda}

\begin{enumerate}[label=\alph*), leftmargin=*, topsep=2pt]
    \item \textbf{Formule} el modelo general que permite colocar piezas de varios
          tipos a la vez maximizando el valor total. Defina con cuidado las nuevas
          variables de decisión (ahora deben distinguir \emph{qué} pieza va en
          \emph{qué} casilla) y las restricciones de ocupación y de no amenaza.
    \item \textbf{Caso 1 --- torres y alfiles.} Resuelva permitiendo únicamente
          torres y alfiles. Reporte el valor total óptimo y cuántas piezas de cada
          tipo se colocan. \textbf{Compare} ese valor con el que se obtendría
          usando solo torres o solo alfiles (con los resultados de la Parte~B).
          ¿Conviene mezclar? ¿Por cuánto?
    \item \textbf{Caso 2 --- torres y caballos.} Resuelva ahora permitiendo torres
          y caballos. Reporte valor total y composición. Compare nuevamente contra
          las estrategias puras. ¿Conviene mezclar en este caso?
    \item \textbf{Caso 3 --- las cuatro piezas.} Resuelva permitiendo los cuatro
          tipos simultáneamente. ¿Qué composición resulta óptima?
    \item \textbf{Análisis.} Los casos anteriores no se comportan igual: en alguno
          de ellos mezclar piezas \emph{mejora} el valor respecto de usar un solo
          tipo, y en otro no. Explique \textbf{por qué}. \emph{Pista:} piense en la
          relación entre el \textbf{valor} de una pieza y cuántas de ellas
          \textbf{caben} en el tablero, es decir, en el ``valor por casilla
          ocupada'' que cada pieza es capaz de sostener.
    \item \textbf{Restricción adicional.} Repita el Caso 3 exigiendo que haya
          \textbf{al menos una pieza de cada tipo} sobre el tablero. ¿Cuánto valor
          se pierde respecto del óptimo libre? Interprete ese costo en términos de
          optimización: ¿qué representa esa diferencia?
\end{enumerate}

\noindent\rule{\linewidth}{0.4pt}

% =============================================================================
\section*{Entregables}
% =============================================================================

\subsection*{1.\quad Informe (PDF)}

Estructura mínima requerida (cada sección comienza en una nueva página):

\begin{enumerate}[leftmargin=*, topsep=2pt]
    \item \textbf{Portada:} nombre del curso, Tarea~3, integrantes, profesor,
          ayudante y fecha.
    \item \textbf{Introducción:} breve explicación del problema, incluyendo la
          \textbf{descripción de las piezas y su comportamiento} en el tablero.
    \item \textbf{Marco teórico:} conceptos utilizados (optimización lineal entera,
          variables binarias, restricciones de incompatibilidad).
    \item \textbf{Metodología y modelamiento:} las formulaciones completas de las
          Partes~A y C.
    \item \textbf{Desarrollo:} breve explicación de la implementación en
          \texttt{gurobipy}.
    \item \textbf{Resultados:} tablas con los máximos por pieza, las
          configuraciones de tablero obtenidas y los valores de la Parte~C.
          \begin{quote}
          \textit{Ojo:} en esta sección van los resultados puros, sin
          interpretación.
          \end{quote}
    \item \textbf{Análisis y conclusiones:} interpretación de los resultados y
          respuestas a las preguntas de análisis. Sea \textbf{específico}; evite
          frases genéricas del tipo ``el modelo funciona''.
    \item \textbf{Referencias} (si aplica).
\end{enumerate}

\begin{ayuda}{Formalidades del informe (se descuentan puntos por incumplimiento)}
\begin{itemize}[topsep=2pt]
    \item Lenguaje técnico y formal, sin errores ortográficos ni gramaticales.
    \item Texto \textbf{justificado}, fuente tamaño \textbf{12} (Times New Roman,
          Arial u otra tipografía estándar).
    \item Cada figura, gráfico o tabla debe llevar \textbf{título y fuente}. Si el
          material es de elaboración propia, debe indicarse explícitamente.
    \item Sean \textbf{breves y directos}. Se prefiere un análisis corto, directo
          y correcto por sobre páginas de relleno.
\end{itemize}
\end{ayuda}

\subsection*{2.\quad Cuaderno Jupyter (.ipynb)}

\begin{itemize}[topsep=2pt]
    \item Implementación \textbf{comentada} de los cuatro modelos de la Parte~B y
          del modelo general de la Parte~C.
    \item Impresión de los resultados y de las configuraciones de tablero
          obtenidas.
    \item Verificación explícita de que las soluciones no contienen amenazas.
    \item Código limpio, ordenado y bien estructurado. \textbf{Un solo archivo de
          código por grupo.}
\end{itemize}

\begin{ayuda}{Puntaje extra: visualización del tablero (+5 pts)}
\vspace{2pt}
Se otorgarán \textbf{5 puntos extra} si presentan una \textbf{visualización gráfica}
de los tableros obtenidos (los cuatro de la Parte~B). Los 5 puntos son por el
conjunto de las visualizaciones, no por cada una. Pueden resolverlo con la
herramienta que prefieran; una alternativa cómoda es la librería \texttt{chess} de
Python, aunque una figura hecha con \texttt{matplotlib} también es perfectamente
válida. La visualización es \textbf{opcional}.
\end{ayuda}

\medskip
\noindent\textbf{Formato de entrega:} comprima el informe y el notebook en un
archivo \texttt{.zip} o \texttt{.rar} con el nombre \texttt{Tarea3\_GrupoN}.
\textbf{No entregue los archivos por separado.}

\noindent\rule{\linewidth}{0.4pt}

% =============================================================================
\section*{Plazos y Evaluación}
% =============================================================================

\begin{itemize}[topsep=2pt]
    \item \textbf{Plazo de entrega:} \textbf{lunes 27 de julio de 2026} a las
          23:59.
    \item \textbf{Grupos:} 2 a 3 integrantes (ni más ni menos).
    \item \textbf{Inscripción:} anótense en los grupos de Canvas
          (Personas~$\to$ Grupos~$\to$ Tarea~3\,--\,Grupo~n).
\end{itemize}

\vspace{0.5em}

\begin{center}
\renewcommand{\arraystretch}{1.2}
\begin{tabular}{lc}
\toprule
\textbf{Criterio} & \textbf{Ponderación} \\
\midrule
Parte A: Formulación de los cuatro modelos      & 40\,\% \\
Parte B: Resolución en Gurobi y análisis        & 30\,\% \\
Parte C: Combinación de piezas por valor        & 30\,\% \\
\midrule
\textit{Visualización de los tableros}          & \textit{+5 pts extra} \\
\textit{Calidad del informe y profundidad del análisis} &
\textit{transversal} \\
\bottomrule
\end{tabular}
\captionof{table}{Criterios de evaluación.}
\end{center}

\vspace{0.3em}
\noindent\textit{La calidad del informe, la profundidad del análisis y la correcta
interpretación de los resultados impactan de forma transversal en todas las
partes.}

\noindent\rule{\linewidth}{0.4pt}

% =============================================================================
\section*{Notas Importantes}
% =============================================================================

\begin{itemize}[topsep=2pt]
    \item \textbf{Deben usar \texttt{gurobipy}} para la implementación.
    \item Recuerden que el \textbf{bloqueo es irrelevante}: basta con impedir que
          dos piezas compartan una misma línea de acción (ver la observación de la
          primera página).
    \item Cada variable debe declararse con su \textbf{dominio y tipo correctos}.
          Identificar el tipo adecuado es parte de la evaluación.
    \item \textbf{Nombren} las restricciones al crearlas: facilita la lectura de la
          solución y la depuración.
    \item Construyan los modelos de forma \textbf{algebraica} (con conjuntos y
          sumatorias sobre ellos); no escriban las 64 casillas una por una.
    \item Pueden apoyarse en herramientas externas para entender el método, pero la
          \textbf{formulación} y el \textbf{análisis} deben ser propios y quedar
          plasmados en el informe.
    \item En el informe y en el código deben figurar los \textbf{nombres de todos
          los integrantes} del grupo.
\end{itemize}

\vspace{1em}
\begin{center}
\textit{Consultas al correo
\href{mailto:vramireza@udd.cl}{\texttt{vramireza@udd.cl}} o al
\texttt{+56\,9\,9912\,1814}
(respondo siempre que puedo, dentro de horarios razonables).}
\end{center}

% =============================================================================
\end{document}
% =============================================================================
